% TOP_PROBLEM: {"id": 30006799, "title": "Serre's Conjecture II on Galois Cohomology of Simply Connected Semisimple Groups", "category_id": 4, "category": {"id": 4, "name": "algebra", "display_name": "Algebra", "description": "Group theory, ring theory, field theory, and algebraic structures.", "slug": "algebra", "order_index": 4, "created_at": "2026-07-31T15:26:25.671Z"}, "status": "open"}
% FIRST_POSTED: 2026-09-27
% !TeX program = lualatex
% Compile twice: lualatex 30006799.tex
% All references and prose are embedded; no BibTeX database or external inputs.
\documentclass[11pt,a4paper]{article}
\usepackage[margin=24mm,headheight=14pt]{geometry}
\usepackage{fontspec}
\setmainfont{Latin Modern Roman}
\setsansfont{Latin Modern Sans}
\setmonofont{Latin Modern Mono}
\usepackage{amsmath,amssymb,mathtools}
\usepackage{unicode-math}
\setmathfont{Latin Modern Math}
\usepackage{microtype}
\usepackage{enumitem,xurl,needspace}
\usepackage[dvipsnames]{xcolor}
\usepackage{hyperref}
\hypersetup{unicode=true,colorlinks=true,linkcolor=MidnightBlue,urlcolor=MidnightBlue,
 pdftitle={Research Notebook: Section 30006799},pdfauthor={Research notebook},bookmarksnumbered=true}
\usepackage{bookmark}
\usepackage{tocloft}
\setlength{\cftsecnumwidth}{3em}
\cftsetpnumwidth{2.5em}
\cftsetrmarg{4em}
\usepackage{titlesec}
\titleformat{\section}{\Large\bfseries\raggedright}{\thesection}{0.7em}{}
\usepackage{fancyhdr}
\pagestyle{fancy}\fancyhf{}
\fancyhead[L]{\small\sffamily Open Problems | Research notebook}
\fancyhead[R]{\small 226 | 27 September 2026}
\fancyfoot[C]{\thepage}
\setlength{\parindent}{0pt}
\setlength{\parskip}{5pt plus 1pt}
\setlength{\emergencystretch}{3em}
\setcounter{tocdepth}{1}
\setcounter{secnumdepth}{1}
\setlist[itemize]{leftmargin=3em,itemsep=5pt,topsep=4pt,parsep=0pt}
\allowdisplaybreaks
\newcommand{\N}{\mathbb{N}}
\newcommand{\Z}{\mathbb{Z}}
\newcommand{\Q}{\mathbb{Q}}
\newcommand{\R}{\mathbb{R}}
\newcommand{\C}{\mathbb{C}}
\newcommand{\F}{\mathbb{F}}
\newcommand{\A}{\mathbb{A}}
\newcommand{\PP}{\mathbb P}
\newcommand{\E}{\mathbb{E}}
\newcommand{\eps}{\varepsilon}
\newcommand{\dd}{\,\mathrm d}
\newcommand{\sref}[2]{\hyperlink{source:#1:#2}{[S#2]}}
\newcommand{\eref}[2]{\hyperlink{extra:#1:#2}{[E#2]}}
% Turn the next switch off for a shorter reading copy. The quotations preserve
% every exactTarget field, including qualifications omitted from a short summary.
\newif\ifshowcatalogue
\showcataloguetrue
\newcommand{\cataloguescope}[1]{\ifshowcatalogue
 \paragraph{Exact scope in the supplied catalogue.}
 \begin{quote}\small #1\end{quote}\fi}

% Independently compilable extraction; original section text retained.
\numberwithin{equation}{section}
\begin{document}
\setcounter{section}{0}
% ================================================================
% problemId: 30006799
% Canonical problem ID: 30006799
% exactTarget SHA-256: ebcf4ee63683559f3765f38aff3bf239bb53600bca902abc2e18afc385f1a10d
\clearpage
\section*{\texorpdfstring{Serre's Conjecture II on Galois Cohomology of Simply Connected Semisimple Groups}{Serre's Conjecture II on Galois Cohomology of Simply Connected Semisimple Groups}}\label{30006799}
\subsection{Definitions and mathematical statement}
Let $K$ be perfect, so every algebraic extension is separable, and write $\Gamma_K=\operatorname{Gal}(K^{\mathrm{sep}}/K)$. Cohomological dimension at most two means vanishing of continuous Galois cohomology in degrees greater than two for all discrete torsion coefficient modules. For a semisimple simply connected algebraic $K$-group $G$, the pointed set $H^1(K,G)$ classifies $G$-torsors. Such a torsor becomes isomorphic to $G$ after a separable extension and is trivial exactly when it has a $K$-point. The assertion is
\[
 H^1(K,G)=\{1\}
\]
for every pair $K,G$ in this scope. Simply connected is the algebraic-group notion of having no nontrivial central isogeny cover, not topological simple-connectedness of a set of real points.
\par\smallskip\textit{Formulation and concept references:} \sref{30006799}{1}.
\cataloguescope{Let K be a perfect field of cohomological dimension at most 2. Prove that every principal homogeneous space under a semisimple simply connected linear algebraic K-group has a rational point, i.e. H\textasciicircum{}1(K,G) = 1 for every such group G.}
\subsection{Short English statement}
Over a perfect field of cohomological dimension at most two, must every torsor under a semisimple simply connected algebraic group have a rational point?
% BEGIN RESEARCH ATTEMPT 226
\subsection{Research attempt}

\paragraph{Current frontier and bibliographic correction.}
The URL in \sref{30006799}{1} points to a different work from the Gille survey
named there: it is Izquierdo--Lucchini Arteche's \emph{Transfer principles
for Galois cohomology and Serre's conjecture II}, whose third version is
dated 9 September 2025. Among its consequences, the characteristic-zero
conjecture implies the positive-characteristic conjecture; it does not
prove the characteristic-zero premise. \eref{30006799}{1}, Introduction and
Section~4.3. Gille's actual survey is separately identified in
\eref{30006799}{2}. Its classical-group and quasi-split exceptional results
are distinguished from the unrestricted assertion.
The March 2026 manuscript of Nguyen extends the formulation to certain
pseudo-reductive groups and proves an equivalence, together with local
and global-field cases, not the general premise. \eref{30006799}{4}, Abstract.
These observations are a source audit, not changes to the original
perfect-field target.

\paragraph{Attack route.}
A torsor can be attacked through an explicit norm equation, through
cohomological invariants, or by transfer from a smaller class of fields.
The first route gives a complete calculation for quaternion norm-one
groups. It reveals exactly why a degree-three invariant is sufficient
there. We then try to transplant that mechanism to a general simply
connected group, keeping the kernel problem visible. Transfer principles
cannot repair a missing invariant-detection theorem merely by changing
characteristic. The lower-rank argument below is a rederivation of a
known subcase, not an assertion of a newly solved family.

\paragraph{Attempt.}
\textbf{1. What the split cases actually use.}
For every field $K$, the exact sequence
\[
 1\longrightarrow\mathrm{SL}_r\longrightarrow\mathrm{GL}_r
       \xrightarrow{\det}\mathbf G_m\longrightarrow1
\]
and Hilbert's theorem~90 for $\mathrm{GL}_r$ identify
$H^1(K,\mathrm{SL}_r)$ with the cokernel of
$\det:\mathrm{GL}_r(K)\to K^*$. The diagonal matrix
$\operatorname{diag}(a,1,\ldots,1)$ shows that this map is surjective.
Thus this family has trivial torsors without any bound on cohomological
dimension. Likewise a split symplectic torsor is a nondegenerate
alternating form on a vector space; choosing a vector, a partner with
pairing one, and their orthogonal complement inductively constructs a
symplectic basis. These are explicit split-group arguments. They do
not trivialize a nonsplit form of the group. The cohomological framework
and its norm-one interpretation are recalled in \eref{30006799}{2}, Sections~1
and~3.

For a central simple algebra $D$, the corresponding exact sequence is
\[
 1\longrightarrow\mathrm{SL}_1(D)\longrightarrow\mathrm{GL}_1(D)
     \xrightarrow{\operatorname{Nrd}}\mathbf G_m\longrightarrow1.
\]
Noncommutative Hilbert~90 gives
\begin{equation}\label{30006799:normtorsors}
 H^1(K,\mathrm{SL}_1(D))=K^*/\operatorname{Nrd}(D^*).
\end{equation}
This is a genuine reduction to surjectivity of a norm, not to the
vanishing of the Brauer class of $D$. A division algebra can remain
nonsplit even when every torsor under its norm-one group is trivial.

\textbf{2. Quaternion norm surjectivity from a dimension gap.}
Assume now $\operatorname{char}K\ne2$ and $\operatorname{cd}(K)\le2$.
Let $D=(u,v)_K$ be a quaternion algebra. Its reduced norm is the
four-dimensional Pfister form
\[
 q=\langle1,-u,-v,uv\rangle.
\]
If $D$ is split, the norm is the determinant on $M_2(K)$ and is already
surjective on $K^*$. Suppose $D$ is a division algebra; then $q$ is
anisotropic. For any $a\in K^*$ consider the eight-dimensional form
\[
 \pi=q\otimes\langle1,-a\rangle=q\perp(-a q).
\]
Let $I$ be the fundamental ideal of the Witt ring. The established
Milnor isomorphisms identify $I^3/I^4$ with
$H^3(K,\Z/2)$, and the class of $\pi$ maps to the symbol
$(u)\cup(v)\cup(a)$. This cohomology group is zero, so $[\pi]\in I^4$.
The Arason--Pfister dimension theorem says that an anisotropic form
with nonzero class in $I^4$ has dimension at least $16$.
The anisotropic part of $\pi$ has dimension at most $8$, so it must
vanish. Both deep inputs, including the characteristic-not-two
convention, are stated in \eref{30006799}{3}, Section~4, Theorem~4.1 and the
Milnor-isomorphism discussion. They are used here as theorems, not
proved by this notebook.

In particular $\pi$ is isotropic. There are $x,y\in D$, not both zero,
with $q(x)=a q(y)$. Anisotropy of $q$ forces $x,y\ne0$; in a division
algebra $y$ is invertible. Multiplicativity of reduced norm gives
\begin{equation}\label{30006799:ratio}
 \operatorname{Nrd}(xy^{-1})=q(x)/q(y)=a.
\end{equation}
Since $a$ was arbitrary, \eqref{30006799:normtorsors} proves this quaternion
subcase. The decisive mechanism is not just that an invariant vanishes:
it is that vanishing moves a specific eight-dimensional form into a
filtration level whose nonzero anisotropic elements require dimension
at least sixteen.

\textbf{3. Attempt to propagate the dimension-gap mechanism.}
For a simply connected absolutely almost simple group in characteristic
zero, the Rost invariant is a pointed map
\[
 r_G:H^1(K,G)\longrightarrow H^3(K,\Q/\Z(2)).
\]
The right side is zero when $\operatorname{cd}(K)\le2$.
Consequently a proof of
\begin{equation}\label{30006799:kernel}
 r_G^{-1}(0)=\{1\}
\end{equation}
for the relevant groups and fields would prove their instances of the
conjecture. For quasi-split exceptional groups of the types covered by
Gille's Theorem~5.2 this is an established route. It is not a theorem
that \eqref{30006799:kernel} holds for arbitrary $G,K$: the same source
describes nontrivial Rost kernels for $E_8$ over fields outside the
present dimension-two scope. \eref{30006799}{2}, Sections~4.2 and~5.2.
They refute an unrestricted invariant-injectivity assertion, not
Serre's dimension-two conjecture.

The quaternion proof suggests looking for a finite filtration on a
geometric object representing the torsor, with two properties: vanishing
of degree-three invariants raises the filtration, and a size bound
forces a terminal filtration class to be trivial. Neither such an
object nor a size bound independent of the torsor has been supplied
for the remaining exceptional groups. Merely introducing further
invariants of degrees $4,5,\ldots$ is not enough: all their ordinary
Galois-cohomology targets also vanish under the given dimension bound.
One still needs a theorem that their joint zero fiber consists only of
the trivial torsor. Defining an invariant with trivial kernel by using
the torsor itself would be circular.

\textbf{4. A restriction--corestriction test.}
A second attempted repair was to trivialize a torsor over finitely many
extensions $K_i/K$ whose degrees have gcd one, then combine these
trivializations using a Bezout identity. For an abelian cohomology group
one often has
\[
 \operatorname{cor}_{K_i/K}\operatorname{res}_{K_i/K}(x)
          =[K_i:K]x,
\]
so such an argument can work. But $H^1(K,G)$ here is a pointed set,
not an abelian group with addition, integer multiples and a universally
available corestriction. The required substitute is the separate
injectivity assertion
\[
 \ker\left(H^1(K,G)\to\prod_i H^1(K_i,G)\right)=\{1\}
        \quad\text{when }\gcd_i[K_i:K]=1.
\]
This is exactly the kind of additional descent question isolated in
\eref{30006799}{2}, Section~4.3. Replacing a rational point by a zero-cycle
of degree one without proving this descent step would weaken the target.
The transfer paper itself distinguishes the zero-cycle and rational-point
versions. \eref{30006799}{1}, Introduction and Section~4.

\paragraph{Stress test.}
Cohomological dimension two does not make all degree-two obstructions
vanish. For example, the Brauer group of a nonarchimedean local field
contains nonzero quaternion classes; they give nontrivial
$\mathrm{PGL}_2$ torsors. This group is not simply connected, so the
example is outside the target and explains why that hypothesis matters.
The distinction and the local-field results are part of
\eref{30006799}{2}, Sections~1--3. Nothing in \eqref{30006799:ratio} asserted that
$D$ itself became split.

The quadratic argument assumes characteristic different from two
explicitly; it is not a proof for every perfect field. The
characteristic-zero Rost calculation likewise avoids importing an
incorrect ordinary-cohomology description of positive-characteristic
primary invariants. The 2025 transfer theorem could carry a full
characteristic-zero solution to the remaining characteristic, but a
proof for quaternion groups alone is not its full premise. Split groups,
inner forms, products and restrictions of scalars must not be identified
without their separate descent arguments. None of the numerical
invariants considered here is being assumed complete.

\needspace{15\baselineskip}
\paragraph{Outcome.}
\textbf{Outcome: Exploratory approach with unresolved gap.}

The attempt supplies the explicit quaternion norm calculation, explains
its filtration/dimension mechanism, and identifies precisely the
invariant-kernel and descent statements needed to extend it. The
calculation uses established deep theorems and belongs to a known
subcase; no new family of Serre-II cases is claimed.

No complete invariant-detection or coprime-degree descent argument has
been obtained in the remaining generality. In particular, the vanishing
of the target of the Rost invariant has not been confused with the
vanishing of its source. Both the all-group and all-perfect-field
quantifiers remain unresolved by this attempt.
% END RESEARCH ATTEMPT 226
\subsection{Sources}
\begin{itemize}\raggedright
\item[\textnormal{[S1]}]\hypertarget{source:30006799:1}{} Philippe Gille, "Serre's conjecture II: a survey," in Quadratic forms, linear algebraic groups, and cohomology, Dev. Math. 18, Springer, New York (2010), 41-56.
\url{https://arxiv.org/abs/2308.00903}.
{\small\textit{Catalogue formulation source.}}
% BEGIN ADDED REFERENCES 226
\item[\textnormal{[E1]}]\hypertarget{extra:30006799:1}{}
D. Izquierdo and G. Lucchini Arteche, \emph{Transfer principles for Galois
cohomology and Serre's conjecture II}, arXiv:2308.00903v3,
9 September 2025, Introduction and Sections~3--4, particularly~4.3.
\url{https://arxiv.org/html/2308.00903v3}.
This is the actual work reached by the original [S1] URL.
\item[\textnormal{[E2]}]\hypertarget{extra:30006799:2}{}
P. Gille, \emph{Serre's conjecture II: a survey}, in \emph{Quadratic Forms,
Linear Algebraic Groups, and Cohomology}, Developments in Mathematics
18, Springer (2010), 41--56. Consulted author manuscript,
Sections~1--5, especially~4.2--4.3 and Theorems~5.1--5.2.
\url{https://www.math.uni-bielefeld.de/lag/man/326.pdf}.
Its historical status discussion is supplemented by [E1] and [E4].
\item[\textnormal{[E3]}]\hypertarget{extra:30006799:3}{}
A. Qu\'eguiner-Mathieu, \emph{Pfister forms in the algebraic and geometric
theory of quadratic forms}, author-hosted exposition, Section~4,
Theorem~4.1 and the Milnor-isomorphism discussion on pp.~6--7.
\url{https://www.math.univ-paris13.fr/~queguin/fichiers/Pfister.pdf}.
The cited inputs are the Arason--Pfister theorem and the established
Milnor conjectures, not conjectural additional hypotheses.
\item[\textnormal{[E4]}]\hypertarget{extra:30006799:4}{}
M. N. T. Nguyen, \emph{Serre conjecture II for pseudo-reductive groups},
arXiv:2603.08061v1, 9 March 2026, Abstract and formulation of the
pseudo-reductive extension. \url{https://arxiv.org/abs/2603.08061}.
All four research references were consulted 27 September 2026.
% END ADDED REFERENCES 226
\end{itemize}

\end{document}
